% ============================================================
% MCM/ICM STARTER TEMPLATE
% ============================================================
% This is a working template for the Mathematical Contest in
% Modeling (MCM) / Interdisciplinary Contest in Modeling (ICM).
%
% IMPORTANT MCM RULES baked into this file:
%   - Papers must be written in ENGLISH.
%   - You must NOT put your name, school, or advisor anywhere in
%     the paper (blind review). Only your team control number
%     goes on the summary sheet.
%   - Papers have a strict PAGE LIMIT (check current-year rules,
%     historically 25 pages including everything but the
%     Report Summary counts toward it in recent years -- always
%     re-check the official rules PDF for the year you compete).
%
% How to use this file:
%   1. Read the comments (%) -- they explain what each block does,
%      what you should change, and common mistakes.
%   2. Delete the example content (marked EXAMPLE) and replace it
%      with your own.
%   3. Compile with pdfLaTeX (or XeLaTeX/LuaLaTeX if you add
%      non-Latin fonts). If you use natbib/bibtex, remember the
%      compile sequence: pdflatex -> bibtex -> pdflatex -> pdflatex.
% ============================================================

\documentclass[12pt, letterpaper]{article}
% [12pt] font size -- MCM has no strict font-size rule, but 12pt
%   body / reasonable margins is the safe, common choice.
% [letterpaper] MCM is a US-run contest -- use letterpaper, not a4paper.
%   WRONG:  \documentclass[a4paper]{article}  (margins won't match
%           the page-limit expectations most teams design around)

% ---------------- PACKAGES ----------------
% Load only what you use. A huge unused preamble is a common
% novice habit -- it doesn't break compilation but it's clutter
% and can occasionally cause package-conflict errors.

\usepackage[utf8]{inputenc}     % lets you type UTF-8 characters directly
\usepackage[T1]{fontenc}        % better font encoding / hyphenation for accented chars
\usepackage{amsmath, amssymb, amsthm}
                                 % amsmath: align, cases, matrices, \text{}
                                 % amssymb: extra symbols like \mathbb{R}, \leq, \geq
                                 % amsthm:  \newtheorem for Theorem/Lemma/Definition boxes
\usepackage{geometry}
\geometry{margin=1in}           % 1-inch margins is the safe default.
                                 % WRONG: shrinking margins to 0.5in to cram in more
                                 % text -- graders notice, and past a certain point
                                 % it looks like you're gaming the page limit.

\usepackage{graphicx}           % \includegraphics{...}
\usepackage{float}              % gives the [H] "put it EXACTLY here" figure placement
\usepackage{caption, subcaption} % nicer captions + side-by-side subfigures
\usepackage{booktabs}           % \toprule \midrule \bottomrule -- professional-looking tables
                                 % WRONG: using \hline everywhere (looks like Excel, not a paper)
\usepackage{multirow}           % merge table cells vertically
\usepackage{array}              % extra column types, e.g. fixed-width columns

\usepackage{algorithm}          % floating "Algorithm N" box
\usepackage{algpseudocode}      % \State, \If, \For, \While, etc. pseudocode commands

\usepackage{listings}           % source-code listings (for your appendix)
\usepackage{xcolor}             % \textcolor{}, and needed by listings for code coloring

\usepackage[
  sort&compress,     % combine consecutive citation numbers like [3-5]
  numbers            % numeric citation style [1], [2], ... (common for MCM)
]{natbib}
% Alternative if your coordinator/advisor prefers author-year style:
%   \usepackage[round]{natbib}   -> gives (Smith, 2020) instead of [1]
% Pick ONE citation style and use it consistently. Mixing \cite styles
% is a common mistake that looks sloppy to judges.

\usepackage{hyperref}           % clickable links/refs -- LOAD THIS LAST
% WRONG: loading hyperref before other packages (e.g. before amsthm,
% algorithm) can break their cross-referencing; hyperref should
% almost always be the last (or second-to-last, before cleveref) package.
\hypersetup{
  colorlinks=true,
  linkcolor=black,   % internal refs (sections/figures) -- black keeps it print-friendly
  citecolor=black,
  urlcolor=blue,
  pdftitle={Your Paper Title},
  pdfauthor={},      % LEAVE BLANK -- do not put your name (blind review)
}

\usepackage{fancyhdr}           % custom headers/footers, e.g. team number on every page
\pagestyle{fancy}
\fancyhf{}
\fancyhead[C]{Team \#XXXXXX}    % <-- replace XXXXXX with your real team control number
\fancyfoot[C]{\thepage}
\renewcommand{\headrulewidth}{0.4pt}

% ---------------- THEOREM-LIKE ENVIRONMENTS ----------------
% Optional -- useful if your model has formal definitions/assumptions
% you want numbered and visually boxed.
\newtheorem{assumption}{Assumption}
\newtheorem{definition}{Definition}

% ============================================================
% DOCUMENT
% ============================================================
\begin{document}

% ---------------- SUMMARY SHEET (Page 1) ----------------
% MCM requires a one-page "Summary" (sometimes called Abstract).
% It is judged heavily -- many teams say it's the single most
% important page, since judges triage papers by it.
% RULES: no name/school here. Team number and problem letter ARE
% allowed (check the current problem statement for the exact
% required header format -- it changes slightly year to year).

\begin{center}
    {\Large\bfseries Your Paper Title Here}\\[6pt]
    % EXAMPLE -- replace with your title. Keep it descriptive,
    % not generic ("A Model for..." is overused and vague).

    \textbf{Team \#XXXXXX}
\end{center}

\begin{abstract}
% EXAMPLE summary structure (fill each part in ~1 short paragraph):
%   1. Restate the problem in your own words (1-2 sentences).
%   2. State your overall approach / key model(s) used.
%   3. State your main results / numbers, not just "we solved it."
%   4. State strengths/limitations briefly.
%   5. One sentence on real-world implication or recommendation.
%
% WRONG: copy-pasting the problem statement here -- judges have
% already read it; restate it, don't repeat it verbatim.
% WRONG: leaving out actual numeric results -- "our model performs
% well" is meaningless without a number or comparison.

This paper addresses \dots We propose a \dots model combining
\dots Our results show that \dots We validate the model via
\dots and find it is robust to \dots but sensitive to \dots
Overall, our approach suggests that \dots

\textbf{Keywords:} keyword one; keyword two; keyword three
% 3-5 keywords is typical.
\end{abstract}

\newpage
\tableofcontents
\newpage

% ============================================================
% 1. INTRODUCTION
% ============================================================
\section{Introduction}

\subsection{Problem Background}
% Give context: why does this problem matter in the real world?
% One paragraph, cite real sources if you use them (\citep{...}).

\subsection{Restatement of the Problem}
% Rewrite the problem in your own precise, technical language.
% This is where you convert a fuzzy prompt into something you
% can actually build equations for.

\subsection{Our Work}
% A short roadmap: "In Section 2 we state our assumptions... In
% Section 3 we build Model I... In Section 4 we test..."
% This costs you almost nothing and helps judges navigate a paper
% they're skimming under time pressure.

% ============================================================
% 2. ASSUMPTIONS AND JUSTIFICATIONS
% ============================================================
\section{Assumptions}
% Every assumption needs a ONE-LINE justification. An unjustified
% assumption reads as "we didn't think about this."
% WRONG:
%   \item Air resistance is negligible.
% RIGHT:
%   \item Air resistance is negligible, since the object's terminal
%     velocity in this regime is over two orders of magnitude
%     larger than its actual speed.

\begin{itemize}
    \item \textbf{Assumption 1:} \dots
    \\ \textit{Justification:} \dots
    \item \textbf{Assumption 2:} \dots
    \\ \textit{Justification:} \dots
\end{itemize}

% ============================================================
% 3. NOTATION
% ============================================================
\section{Notation}
% A symbol table. Judges (and you, three days into the contest)
% will thank you. Put it early, before your model, not buried
% in an appendix.

\begin{table}[H]
    \centering
    \caption{Notation used in this paper}
    \begin{tabular}{cl}
        \toprule
        \textbf{Symbol} & \textbf{Meaning} \\
        \midrule
        $x$        & description of $x$ \\
        $\theta$   & description of $\theta$ \\
        $t$        & time (units: \dots) \\
        \bottomrule
    \end{tabular}
    \label{tab:notation}
\end{table}
% ALWAYS give units where relevant -- a common judge complaint is
% "no units anywhere in the model."

% ============================================================
% 4. MODEL
% ============================================================
\section{Model Development}

\subsection{Overview}
% One paragraph: what kind of model (ODE, optimization, ML,
% simulation, graph-theoretic, ...) and why that choice fits
% the problem, before diving into equations.

\subsection{Model Formulation}

% ---- Single numbered equation ----
Consider the relationship between $x$ and $t$ given by
\begin{equation}
    \frac{dx}{dt} = r x \left(1 - \frac{x}{K}\right)
    \label{eq:logistic}
\end{equation}
where $r$ is the growth rate and $K$ is the carrying capacity.
% \label + \eqref lets you refer back: "as shown in
% Equation~\eqref{eq:logistic}". Always \label equations you cite
% later -- don't hardcode "Equation (3)" as literal text, since
% renumbering will silently break it.

% ---- Multi-line aligned equations ----
% Use `align` (not several separate `equation`s) when you want
% equals signs lined up, e.g. deriving a result step by step:
\begin{align}
    f(x) &= a x^2 + bx + c \\
         &= a\left(x + \frac{b}{2a}\right)^2 + \left(c - \frac{b^2}{4a}\right)
\end{align}
% WRONG: using multiple $$ ... $$ blocks for this -- $$ is
% deprecated LaTeX2e style; use \[ \] or, better, a numbered
% environment like equation/align so you can reference it.

% ---- Piecewise definition ----
\begin{equation}
    g(x) =
    \begin{cases}
        1 & \text{if } x \geq 0, \\
        0 & \text{if } x < 0.
    \end{cases}
\end{equation}

% ---- Matrix ----
\begin{equation}
    A =
    \begin{pmatrix}
        a_{11} & a_{12} \\
        a_{21} & a_{22}
    \end{pmatrix}
\end{equation}
% bmatrix gives [ ], pmatrix gives ( ), vmatrix gives | | (determinant-style).

\subsection{Algorithm}
% If your model involves a procedure (simulation, optimization
% loop, etc.), a pseudocode box is often clearer than prose.
\begin{algorithm}[H]
    \caption{Example iterative update}
    \label{alg:example}
    \begin{algorithmic}[1]
        \State Initialize $x_0$, tolerance $\epsilon$, max iterations $N$
        \For{$k = 0, 1, \dots, N$}
            \State $x_{k+1} \gets x_k - \eta \nabla f(x_k)$
            \If{$|x_{k+1} - x_k| < \epsilon$}
                \State \textbf{break}
            \EndIf
        \EndFor
        \State \Return $x_{k+1}$
    \end{algorithmic}
\end{algorithm}

% ============================================================
% 5. FIGURES AND TABLES
% ============================================================
\section{Results}

% ---- Single figure ----
% \begin{figure}[H]
%     \centering
%     \includegraphics[width=0.7\textwidth]{figures/my_plot.png}
%     \caption{Description of what this figure shows and why it matters.}
%     \label{fig:my-plot}
% \end{figure}
% Commented out because the image file doesn't exist yet --
% uncomment once you have figures/my_plot.png. Compiling with a
% missing image path is a very common first error.
%
% WRONG: \caption{Figure 1.} -- a caption should describe the
% content, not just label it; LaTeX numbers it for you already.

% ---- Two subfigures side by side ----
% \begin{figure}[H]
%     \centering
%     \begin{subfigure}[b]{0.45\textwidth}
%         \includegraphics[width=\textwidth]{figures/left.png}
%         \caption{Left condition}
%     \end{subfigure}
%     \hfill
%     \begin{subfigure}[b]{0.45\textwidth}
%         \includegraphics[width=\textwidth]{figures/right.png}
%         \caption{Right condition}
%     \end{subfigure}
%     \caption{Overall caption comparing both panels.}
%     \label{fig:comparison}
% \end{figure}

% ---- Table with booktabs (the "professional" table style) ----
\begin{table}[H]
    \centering
    \caption{Example results table}
    \begin{tabular}{lcc}
        \toprule
        \textbf{Scenario} & \textbf{Metric A} & \textbf{Metric B} \\
        \midrule
        Baseline    & 0.42 & 1.30 \\
        Variant 1   & 0.51 & 1.12 \\
        Variant 2   & 0.39 & 1.45 \\
        \bottomrule
    \end{tabular}
    \label{tab:results}
\end{table}
% WRONG table habits to avoid:
%   - vertical rules (|) between every column -- visually noisy
%   - \hline above and below every row -- use booktabs rules instead
%   - forgetting \label / \caption -- every float should have both

% ============================================================
% 6. SENSITIVITY ANALYSIS
% ============================================================
\section{Sensitivity Analysis}
% Judges specifically look for this section. Pick 1-2 key
% parameters, vary them, and show/discuss how much the output
% changes. "Our model is robust because..." backed by a number
% or plot is much stronger than asserting robustness with no test.

% ============================================================
% 7. STRENGTHS AND WEAKNESSES
% ============================================================
\section{Strengths and Weaknesses}

\subsection{Strengths}
\begin{itemize}
    \item \dots
\end{itemize}

\subsection{Weaknesses}
\begin{itemize}
    \item \dots
\end{itemize}
% Being honest here is rewarded, not punished -- a model with no
% acknowledged weaknesses reads as either untested or unaware.

% ============================================================
% 8. CONCLUSION
% ============================================================
\section{Conclusion}
% Summarize findings and, if the problem calls for it, give a
% concrete recommendation ("we recommend policy X because...").

% ============================================================
% REFERENCES
% ============================================================
% Two common approaches -- pick ONE:
%
% (A) Manual bibliography (fine for MCM's typically small
%     reference lists, and avoids a BibTeX build step):
\begin{thebibliography}{9}
    \bibitem{smith2020}
    A. Smith, ``Title of the paper,'' \textit{Journal Name}, vol. 1,
    pp. 1--10, 2020.

    \bibitem{website2023}
    Organization Name, ``Title of the page,'' 2023. [Online].
    Available: \url{https://example.com}. [Accessed: Aug. 22, 2026].
\end{thebibliography}
% Cite these in text with \cite{smith2020}.
%
% (B) BibTeX (better for many references; needs a references.bib file):
%   \bibliographystyle{plainnat}
%   \bibliography{references}
% and compile pdflatex -> bibtex template -> pdflatex -> pdflatex.
%
% WRONG: mixing (A) and (B) in the same document, or citing a
% source in text that never appears in the bibliography (and
% vice versa) -- both are easy to do by accident, so check before
% submitting.

% ============================================================
% APPENDIX
% ============================================================
\appendix
\section{Code}
% If you used code (Python/MATLAB/R) to generate results, include
% the key parts here so judges can verify your methodology.
\begin{lstlisting}[language=Python, basicstyle=\small\ttfamily, breaklines=true]
import numpy as np

def logistic_growth(x0, r, K, steps):
    x = x0
    trajectory = [x]
    for _ in range(steps):
        x = x + r * x * (1 - x / K)
        trajectory.append(x)
    return np.array(trajectory)
\end{lstlisting}
% \lstinputlisting[language=Python]{code/my_script.py}
% is often better than pasting code inline -- it pulls directly
% from your actual script file, so you can't accidentally paste
% a stale/edited-after-the-fact version.

\end{document}
